\documentclass[11pt,a4paper]{article}
\usepackage[margin=23mm]{geometry}
\usepackage{fontspec}\setmainfont{Latin Modern Roman}
\usepackage{amsmath,amssymb,amsthm,xurl,hyperref}
\setlength{\parindent}{0pt}\setlength{\parskip}{6pt}
\setlength{\emergencystretch}{3em}
\newcommand{\code}[1]{\nolinkurl{#1}}
\newcommand{\E}{\mathbb E}
\title{S05 / Kolekcja obstructions\\\large Materiał roboczy do ewentualnego Aneksu B}
\author{Notatka robocza dla Niirmata}\date{10 października 2026}
\begin{document}\maketitle
\textbf{WARUNKOWY / DEVELOPMENT.} Kolekcja zachowuje ograniczenia znanych
wyników. Nie jest zmianą paperu, jego publikacją ani deklaracją bezpieczeństwa
lub niebezpieczeństwa FT1536. Manifest \code{INDEX.json} identyfikuje
zakres, status dowodu, źródła, receipty oraz SHA256 każdego wpisu.
Nowe wpisy dodajemy bez przepisywania historycznych źródeł.

\section{O1: DyadicObstruction --- dokładna jednostajność}
\textbf{Teza.} Żaden deterministyczny algorytm zużywający najwyżej $b$
niezależnych uczciwych bitów nie może zwracać dokładnie jednostajnego
współczynnika w $\mathbb Z/18433\mathbb Z$, dla skończonego $b$.

\textbf{Dowód.} Dopełniamy każdą krótszą taśmę do długości $b$.
Masa dowolnego zdarzenia ma wtedy postać $k/2^b$, $k\in\mathbb N$.
Jednostajność wymaga $18433k=2^b$, co jest niemożliwe: nieparzysta liczba
18433 większa od 1 nie dzieli potęgi 2. Primalność nie jest potrzebna
do argumentu tekstowego, choć szkic Lean używa jej jako wygodnego lematu.

\textbf{Dokładny zakres.} Stały skończony limit bitów i dokładna
jednostajność współczynnika. Wyklucza wspólną instancję obecnego typu
taśmy \code{Games.Sampler} oraz \code{HacGlue.UniformChallengeAt}.
Nie obala twierdzenia z tymi warunkowymi przesłankami.
Nie wyklucza aproksymacji ani losowania z nieograniczoną długością taśmy.
Idealna wyrocznia jednostajna jest innym obiektem niż taki sampler.

\textbf{Stan dowodu.} Tekstowy argument; historyczny szkic
\code{DyadicObstruction.lean} istnieje, lecz nie ma potwierdzonego
uruchomienia jego kernela. Historyczny batch zatrzymał się na timeout
wcześniejszego modułu StateErasure; drugi timeout dotyczył samego importu.
Nie należy przedstawiać tych timeoutów jako błędu twierdzenia ani jako PASS.

\textbf{Droga obejścia.} Dla $N=2^b=aq+r$, wynik słowa modulo $q$ ma
dokładny kierunkowy moment względem jednostajnego współczynnika
$1+r(q-r)/N^2$. Przybliżenie musi być opłacone w nowym certyfikacie,
zamiast zmiany znaczenia słowa „jednostajny”.

\section{O2: duży moment przed normą przy pominięciu ważenia marginalu}
\textbf{Konstrukcja objęta wynikiem.} Sampler najpierw losuje $z_2$ z $U=w/G$,
a potem $z_1$ z dokładnego Gaussa w reszcie $c-hz_2$, na skończonym boxie.
Uczciwy Gaussian na całym włóknie ma inne marginalne prawo $z_2$.
Dokładny moment tej konstrukcji wynosi
\[
 \sum S_{h,c}^2/P_{h,c}
 =\E_U[Z_1(c-hz_2)]\E_U[1/Z_1(c-hz_2)].
\]
Przy h=0 czynnik wynosi 1. Przy parametrach FT1536, h=1,c=0,
NORMALIZER-005 dowodzi dolnej granicy $>4$ dla każdego niezależnego bloku,
więc dla 768 bloków moment jest $>2^{1536}$.

\textbf{Mechanizm dowodu.} W kwadracie $[-3072,3072]^2$ masa reszty
jest zdominowana przez centralny lift: każdy inny lift ma energię większą
o co najmniej $q^2-4q\cdot3072>0$. Dolna granica skończonej sumy
odwrotności oraz certyfikowane normalizatory theta dają wynik.
Źródło i pełne końce przedziałów: \code{NORMALIZER_005.tex},
\code{FIBER_NORMALIZER_005_CORRECTED.sage},
\code{NORMALIZER_005_CORRECTED_CERTIFICATE.json}.
Erratum \code{NORMALIZER_005_ERRATUM_001.tex} zastępuje wcześniejsze
numeryczne świadectwa po wykryciu niepoprawnej konwersji końców MPFR;
zakres i dolna granica twierdzenia pozostają takie same po ponownym rachunku.

\textbf{Dokładny zakres.} Warunkowo przy h=1,c=0, \emph{przed normą}.
To nie jest dolna granica momentu po normie, cap16, emisji ani uśrednieniu
po wyzwaniu. Nie dowodzi słabości bezpieczeństwa FT1536 ani niemożliwości
innego publicznego samplera. Status: textual, rygorystyczny Sage512/768,
bez kernela i niezależnego odbioru.

\textbf{Droga obejścia.} Losować marginal proporcjonalny do
$w(z_2)Z_1(c-hz_2)$; przedziały normalizatora z 005 umożliwiają względną
kontrolę tych wag. Nie wystarcza większa precyzja dla błędnego marginalu.

\section{O3: koszt odrzucania całych wektorów dla ważonego marginalu}
\textbf{Konstrukcja objęta wynikiem.} Dokładna propozycja $U=w/G$ i moneta
$Z_1(c-hz_2)/M$, gdzie jedna stała $M$ dominuje wszystkie wagi reszt.
Dla h=1,c=0 każdy taki $M\ge Z_1(0)\ge1$ i certyfikat MARGINAL-006 daje
\[
 \alpha<2^{-767},\qquad
 \Pr[\text{akceptacja w co najwyżej }2^{128}\text{ propozycjach}]<2^{-639}.
\]
Ten wynik pochodzi z faktoryzacji rzeczywistego włókna i rachunku balls
normalizatorów, nie z pomiaru czasu ani eksperymentu o małej liczbie prób.

\textbf{Dokładny zakres.} Przeszkoda kosztowa dla odrzucania całego wektora
z tej propozycji. Nie wyklucza propozycji adaptowanej, faktoryzacji blokowej
ani innego samplera. Nie jest nowym dolnym oszacowaniem końcowego momentu.
Przy h=1 faktoryzacja blokowa usuwa ten konkretny iloczyn prawdopodobieństw
akceptacji; jej pełny skończonobitowy certyfikat pozostaje osobnym zadaniem.

\textbf{Stan dowodu.} Textual, z certyfikatem Sage512/768.
Nowy ogólny lemat capu zapisuje pełną porażkę jako
$J=(1-\delta)Q+\delta\delta_{\mathtt{none}}$ i rozlicza jej moment.
Mała $\delta$ sama nie wystarcza: potrzebna jest także masa
$P(\mathtt{none})$ w mianowniku. Nie przeniesiono certyfikatu h=0
na tę konstrukcję.

\section{Zasady późniejszego użycia w Aneksie B}
Przy każdym wpisie zachować kwantyfikatory, etap prawa, status kernela,
kierunek ilorazu $J/P$ i non-claims. Przeszkody dotyczą określonych typów
lub konstrukcji, nie QROM jako całości. Nie zastępować ich nagłówkiem
„brak bezpieczeństwa”. Kernelizacja h=0 i właściwego lematu capów pozostaje
na żywej liście; szkic Lean O1 nie służy jako dowód, że ten obowiązek wykonano.
\end{document}
